\documentclass[10pt,a4paper]{amsart}
\usepackage[margin=19mm]{geometry}
\usepackage{amsmath,amssymb,mathtools}
\usepackage[scr=rsfs,cal=euler]{mathalfa}
\usepackage{xfrac}
\usepackage[unicode,hidelinks,bookmarks=false]{hyperref}
\newcommand{\ie}{i.e.\ }
\newcommand{\cf}{cf.\ }
\newcommand{\resp}{respectively\ }
\newcommand{\Subsection}{\subsection}
\providecommand{\nobreakdash}{\nobreak\hbox{-}}
\setlength{\emergencystretch}{2em}
\multlinegap0pt
\usepackage{amssymb}
\usepackage{float}
\newtheorem{corollary}{Corollary}
\newtheorem{theorem}{Theorem}
\title{Connectivity for slice-projections of connected polymatroids}
\author{Xiaxia Guan$^{a}$,~~Xian'an Jin$^{b,c}$}
\date{Source: arXiv:2606.22819v1 (CC BY 4.0)}
\begin{document}
\maketitle

\noindent\emph{Source and licence.} Connectivity for slice-projections of connected polymatroids. Version arXiv:2606.22819v1 (CC BY 4.0), distributed by arXiv under the Creative Commons Attribution 4.0 International licence, http://creativecommons.org/licenses/by/4.0/, as recorded in the arXiv licence metadata for this exact version and shown on the abstract page https://arxiv.org/abs/2606.22819v1. Full licence notice: \texttt{licenses/arxiv-2606.22819-CC-BY-4.0.txt}.\par\smallskip

\noindent\textit{Editorial scope.} Counted unit: the article's theorem dual-equal1 (source0.tex lines 415--421). The article's complete original proof of it is delivered with it (source0.tex lines 424--441, one \texttt{proof} environment of 18 lines). No mathematics is written, repaired or re-derived: the statement and the proof are the article's own lines, in the article's own notation, and no retained line is substituted or re-flowed.  Retained uncounted as the source-local context the proof needs: lines 149--151; lines 178--184; lines 406--408.  Nothing was omitted from any retained span, so the package carries no omission record.  No retained line contains a citation, so the wrapper carries no bibliography and no bibliography entry.  No retained line contains a figure environment, an image-inclusion command or drawing code, so no image is delivered and the package declares no asset dependency.  Editorial formatting only, in addition: an AMS \texttt{amsart} preamble that restates the article's own declarations and macros used by the retained lines, namely the environments \texttt{corollary}, \texttt{theorem} on separate counters, together with the standard packages those lines need.  The retained environments print in this document as Corollary 1, Theorem 1 in order of appearance, whereas the article prints the same items with the numbering of its own full document.  The complete native source of the article as distributed in the arXiv e-print for this exact version is bundled unchanged as \texttt{sources/203383/source0.tex}.
\begin{equation}\label{dual-rank}
f^*(I) = f(E-I) + \sum_{i \in I} f(i) - f(E) ~\text{for any}~I\subseteq E.
\end{equation}

 \begin{corollary}\label{dual-equal}
Let $P$ be a polymatroid on $E$ with dual polymatroid $P^*$.  Then for any $e\in E$,
\begin{enumerate}
\item [(i)] $P/e$ and $P^*\setminus e$ are equally connected;
\item [(ii)] $P\setminus e$ and $P^*/e$ are equally connected.
\end{enumerate}
\end{corollary}

\begin{equation}\label{dual-rank1}
r^*(I) = r(E-I) + |I| - r(E) ~\text{for any}~I\subseteq E.
\end{equation}

\begin{theorem}\label{dual-equal1}
For a matroid $M$ on $E$ with dual matroid $M^*$.  Then for any $e\in E$,
\begin{enumerate}
\item [(1)] $M/e$ and $M^*-e$ are equally connected;
\item [(2)] $M-e$ and $M^*/e$ are equally connected.
\end{enumerate}
\end{theorem}

\begin{proof}
Note that the connectedness of a matroid and its corresponding polymatroid are identical.
Let $P=P(M)$ and $P(M^*)$ be the polymatroids corresponding to $M$ and $M^*$, respectively. Then $P-e$ and $P/e$ are the polymatroids corresponding to $M-e$ and $M/e$, respectively.

Let $r$ and $r^*$ be the rank functions of matroids $M$ and $M^*$, respectively. If $r(e')\neq 0$ for any $e'\in E$, then $P^*=P(M^*)$ by Equalities (\ref{dual-rank}) and (\ref{dual-rank1}). This implies that $P^*-e$ and $P^*/e$ are the polymatroids corresponding to $M^*-e$ and $M^*/e$, respectively. By Corollary \ref{dual-equal}, the conclusions hold.

If $r(e')= 0$ for some $e'\in E-e$, then $r^*(E-e')=r^*(E)-1$. Hence, $M/e$, $M^*-e$,  $M-e$ and $M^*/e$ are disconnected. Let $f$ and $f^*$ be the rank functions of polymatroids $P$ and $P^*$, respectively. Clearly, $f(e')=f^*(e')=0$, so, $|T^*_{e'}|=|T_{e'}|=1$. Hence, $P/e$, $P^*-e$,  $P- e$ and $P^*/e$ are also disconnected. In this case, the conclusions hold.


We now assume that $r(e)= 0$ and $r(e')\neq 0$ for any $e'\in E-e$. Then 	
\begin{equation*}\label{Tutte-DC}
				f^*(I)=\left\{\begin{array}{ll}
					r^*(I),&\text{ if } e\notin I,\\
                    r^*(I)-1,&\text{ if } e\in I.
				\end{array}\right.
			\end{equation*}
This implies that $P^*/e=P(M^*)/e$ and $P^*- e=P(M^*)-e$ by the definitions of contraction and deletion. Hence, the conclusions also hold by Corollary \ref{dual-equal}.
\end{proof}

\end{document}
